\ifdefined\gkver\else\def\gkver{student}\fi
\documentclass[\gkver]{gaokaozhenti}
\begin{document}

\chapter{2011年广东理综}
%% paperType: 理综物理部分

\begin{enumerate}
\item
%% number: 13
%% paperName: 2011年广东理综第 13 题 4 分
%% typeId: 1
%% type: 单选题
%% chapter: 气体性质
%% point: 分子力
%% method: 无
%% score: 4
%% degree: 650
%% duplicateId: 0
%% body:
如图所示，两个接触面平滑的铅柱压紧后悬挂起来，下面的铅柱不脱落，主要原因是\xzanswer{D}

\onepicture[w=3.2cm]{figs/13.png}

\fourchoices[answer=D]
{铅分子做无规则热运动}
{铅柱受到大气压力作用}
{铅柱间存在万有引力作用}
{铅柱间存在分子引力作用}

%% answer: D
%% memo:
\memoanswer{
由于铅柱较软，由于接触面平滑后，用力压紧，使得铅分子间的距离小到分子力起作用的距离，分子引力的作用使铅柱在钩码的牵引下未分开，D 正确。
}

\item
%% number: 14
%% paperName: 2011年广东理综第 14 题 4 分
%% typeId: 1
%% type: 单选题
%% chapter: 气体性质
%% point: 热力学第一定律
%% method: 无
%% score: 4
%% degree: 650
%% duplicateId: 0
%% body:
如图所示为某种椅子与其升降部分的结构示意图，M、N 两筒间密闭了一定质量的气体，M 可沿 N 的内壁上下滑动。设筒内气体不与外界发生热交换，在 M 向下滑动的过程中\xzanswer{A}

\onepicture[w=5.2cm]{figs/14.png}

\fourchoices[answer=A]
{外界对气体做功，气体内能增大}
{外界对气体做功，气体内能减小}
{气体对外界做功，气体内能增大}
{气体对外界做功，气体内能减小}

%% answer: A
%% memo:
\memoanswer{
筒内气体不与外界发生热交换，当气体体积变小时，则外界对气体做功，外界对气体做功使气体的内能增大。A 正确。
}

\item
%% number: 15
%% paperName: 2011年广东理综第 15 题 4 分
%% typeId: 1
%% type: 单选题
%% chapter: 电磁感应
%% point: 法拉第电磁感应定律
%% method: 无
%% score: 4
%% degree: 650
%% duplicateId: 0
%% body:
将闭合多匝线圈置于仅随时间变化的磁场中，线圈平面与磁场方向垂直，关于线圈中产生的感应电动势和感应电流，下列表述正确的是\xzanswer{C}

\fourchoices[answer=C]
{感应电动势的大小与线圈的匝数无关}
{穿过线圈的磁通量越大，感应电动势越大}
{穿过线圈的磁通量变化越快，感应电动势越大}
{感应电流产生的磁场方向与原磁场方向始终相同}

%% answer: C
%% memo:
\memoanswer{
由法拉第电磁感应定律，$E = n \frac{\Delta \varphi}{\Delta t} = n \frac{S \Delta B}{\Delta t}$，选项 A 错误。穿过线圈的磁通量越大，并不代表穿过线圈的磁通量变化率大，选项 B 错误，C 正确。由楞次定律感应电流的磁场总是阻碍产生感应电流的磁通量的变化，感应电流的磁场方向与原磁场方向有时相同，有时相反。选项 D 错误。
}

\item
%% number: 16
%% paperName: 2011年广东理综第 16 题 4 分
%% typeId: 1
%% type: 单选题
%% chapter: 力物体的平衡
%% point: 三力平衡
%% method: 无
%% score: 4
%% degree: 650
%% duplicateId: 0
%% body:
如图所示，水平面上，橡皮绳一端固定，另一端连接两根弹簧，连接点 P 在 $F_{1}$、$F_{2}$ 和 $F_{3}$ 三力作用下保持静止，下列判断正确的是\xzanswer{B}

\onepicture[w=3.8cm]{figs/16.png}

\fourchoices[answer=B]
{$F_{1}>F_{2}>F_{3}$}
{$F_{3}>F_{1}>F_{2}$}
{$F_{2}>F_{1}>F_{3}$}
{$F_{3}>F_{2}>F_{1}$}

%% answer: B
%% memo:
\memoanswer{
由于在 $F_{1}$、$F_{2}$ 和 $F_{3}$ 三力作用下保持静止，合力为零，现力 $F_{1}$ 与 $F_{2}$ 垂直，根据力的平行四边形定则由角度及几何关系可得：$F_{3}>F_{1}>F_{2}$，B 正确。
}

\item
%% number: 17
%% paperName: 2011年广东理综第 17 题 6 分
%% typeId: 2
%% type: 多选题
%% chapter: 曲线运动
%% point: 平抛运动
%% method: 无
%% score: 6
%% degree: 450
%% duplicateId: 0
%% body:
如图所示，在网球的网前截击练习中，若练习者在球网正上方距地面 $H$ 处，将球以速度 $v$ 沿垂直球网的方向击出，球刚好落在底线上，已知底线到网的距离为 $L$，重力加速度取 $g$，将球的运动视作平抛运动，下列表述正确的是\xzanswer{AB}

\onepicture[w=4.2cm]{figs/17.png}

\fourchoices[answer=AB]
{球的速度 $v$ 等于 $L\sqrt{\frac{g}{2H}}$}
{球从击出至落地所用时间为 $\sqrt{\frac{2H}{g}}$}
{球从击球点至落地点的位移等于 $L$}
{球从击球点至落地点的位移与球的质量有关}

%% answer: AB
%% memo:
\memoanswer{
球做平抛运动，平抛运动是水平方向上的匀速直线运动和竖直方向上的自由落体运动的合运动，球的初速度 $v = L\sqrt{\frac{g}{2H}}$，A 正确。球从击出到落地的时间 $t = \sqrt{\frac{2H}{g}}$，B 正确。球从击球点至落地点的位移等于 $\sqrt{L^2 + H^2}$，与球的质量无关，选项 C、D 错误。
}

\item
%% number: 18
%% paperName: 2011年广东理综第 18 题 6 分
%% typeId: 2
%% type: 多选题
%% chapter: 光学
%% point: 光电效应
%% method: 无
%% score: 6
%% degree: 650
%% duplicateId: 0
%% body:
光电效应实验中，下列表述正确的是\xzanswer{CD}

\fourchoices[answer=CD]
{光照时间越长光电流越大}
{入射光足够强就可以有光电流}
{遏止电压与入射光的频率有关}
{入射光频率大于极限频率才能产生光电子}

%% answer: CD
%% memo:
\memoanswer{
光电流的大小与入射光的强度相关，A 错误。产生光电效应的条件是：入射光的频率大于或等于被照射材料的极限频率。入射光的频率达不到极限频率，增加照射光的强度是不能产生光电流的，所以 B 错误，C、D 正确。
}

\item
%% number: 19
%% paperName: 2011年广东理综第 19 题 6 分
%% typeId: 2
%% type: 多选题
%% chapter: 交流电
%% point: 变压器
%% method: 无
%% score: 6
%% degree: 450
%% duplicateId: 0
%% body:
图（a）左侧的调压装置可视为理想变压器，负载电路中 $R=55\UO$，A、V 为理想电流表和电压表。若原线圈接入如图（b）所示的正弦交变电压，电压表的示数为 $110\UV$，下列表述正确的是\xzanswer{AC}

\twopicture[numA=（a）, numB=（b）, labelprefix={}, wA=3.4cm, wB=5cm]
{figs/19a.png}
{figs/19b.png}

\fourchoices[answer=AC]
{电流表的示数为 $2\UA$}
{原、副线圈匝数比为 $1:2$}
{电压表的示数为电压的有效值}
{原线圈中交变电压的频率为 $100\UHz$}

%% answer: AC
%% memo:
\memoanswer{
交流电压表和交流电流表的示数都是有效值，电流表的示数 $I = \frac{U}{R} = \frac{110}{55}\UA = 2\UA$，A、C 正确。由图（b）可知原线圈电压有效值为 $220\UV$，周期 $0.02\Us$，则可得原、副线圈匝数比为 $2:1$，交变电压的频率为 $50\UHz$，B、D 错误。
}

\item
%% number: 20
%% paperName: 2011年广东理综第 20 题 6 分
%% typeId: 2
%% type: 多选题
%% chapter: 曲线运动
%% point: 人造地球卫星
%% method: 无
%% score: 6
%% degree: 450
%% duplicateId: 0
%% body:
已知地球质量为 $M$，半径为 $R$，自转周期为 $T$，地球同步卫星质量为 $m$，引力常量为 $G$，有关同步卫星，下列表述正确的是\xzanswer{BD}

\fourchoices[answer=BD]
{卫星距地面的高度为 $\sqrt[3]{\frac{GMT^2}{4\pi^2}}$}
{卫星的运行速度小于第一宇宙速度}
{卫星运行时受到的向心力大小为 $G\frac{Mm}{R^2}$}
{卫星运行的向心加速度小于地球表面的重力加速度}

%% answer: BD
%% memo:
\memoanswer{
卫星距地面的高度为 $\sqrt[3]{\frac{GMT^2}{4\pi^2}} - R$，A 错误。第一宇宙速度是最小的发射卫星的速度，卫星最大的环绕速度，B 正确。同步卫星距地面有一定的高度 $h$，受到的向心力大小为 $G\frac{Mm}{(R + h)^2}$，C 错误。卫星运行的向心加速度为 $\frac{GM}{(R + h)^2}$，地球表面的重力加速度为 $\frac{GM}{R^2}$，D 正确。
}

\item
%% number: 21
%% paperName: 2011年广东理综第 21 题 6 分
%% typeId: 2
%% type: 多选题
%% chapter: 电场
%% point: 带电粒子的直线运动
%% method: 无
%% score: 6
%% degree: 450
%% duplicateId: 0
%% body:
如图所示为静电除尘器除尘机理的示意图。尘埃在电场中通过某种机制带电，在电场力作用下向集尘极迁移并沉积，以达到除尘目的。下列表述正确的是\xzanswer{BD}

\onepicture[w=4.3cm]{figs/21.png}

\fourchoices[answer=BD]
{到达集尘极的尘埃带正电荷}
{电场方向由集尘板指向放电极}
{带电尘埃所受电场力的方向与电场方向相同}
{同一位置带电荷量越多的尘埃所受电场力越大}

%% answer: BD
%% memo:
\memoanswer{
由于集电极与电源正极连接，电场方向由集尘板指向放电极，B 正确。而尘埃在电场力作用下向集尘极迁移并沉积，说明尘埃带负电荷，A 错误。负电荷在电场中受电场力的方向与电场方向相反。C 错误，根据 $F=qE$ 可得，D 正确。
}

\item
%% number: 34
%% paperName: 2011年广东理综第 34 题 12 分
%% typeId: 4
%% type: 实验题
%% chapter: 电路
%% point: 伏安法测电阻
%% method: 无
%% score: 12
%% degree: 450
%% duplicateId: 0
%% body:
\begin{enumerate}
	\item
	（6 分）如图 1 所示是“研究匀变速直线运动”实验中获得的一条纸带，O、A、B、C、D 和 E 为纸带上六个计数点。加速度大小用 $a$ 表示。

	\twopicture[numA=1, numB=2, wA=8cm, wB=6cm]
	{figs/34_1a.png}
	{figs/34_1b.png}

	\begin{steps}
		\item
		OD 间的距离为\tkanswer[1.5cm]{1.20}$\Ucm$。

		\item
		图 2 是根据实验数据绘出的 $s-t^2$ 图线（$s$ 为各计数点至同一起点的距离），斜率表示\tkanswer[2.6cm]{加速度一半}，其大小为\tkanswer[1.8cm]{0.933}$\Umsq$（保留三位有效数字）。
	\end{steps}

	\item
	（12 分）在“描绘小电珠的伏安特性曲线”实验中，所用器材有：小电珠（$2.5\UV$、$0.6\UW$），滑动变阻器，多用电表，电流表，学生电源，开关，导线若干。

	\threepicture[numA=3, numB=4, numC=5, wA=5.2cm, wB=5.4cm, hC=6.5cm]
	{figs/34_2a.png}
	{figs/34_2b.png}
	{figs/34_2c.png}

	\begin{steps}
		\item
		粗测小电珠的电阻，应选择多用电表\tkanswer[1.6cm]{$\times1$}倍率的电阻挡（请填空“$\times1$”“$\times10$”或“$\times100$”）；调零后，将表笔分别与小电珠的两极连接，示数如图 3 所示，结果为\tkanswer[1.4cm]{8}$\UO$。

		\item
		实验中使用多用电表测量电压，请根据实验原理图 4 完成实物图 5 中的连线。

		\item
		开关闭合前，应将滑动变阻器的滑片 P 置于\tkanswer[1.2cm]{a}端。为使小电珠亮度增加，P 应由中点向\tkanswer[1.2cm]{b}端滑动。

		\item
		下表为电压等间隔变化测得的数据。为了获得更准确的实验图像，必须在相邻数据点\tkanswer[1.2cm]{ab}间多测几组数据（请填写“ab”“bc”“cd”“de”“ef”）。

		\begin{center}
		\begin{tabular}{|c|c|c|c|c|c|c|}
		\hline
		数据点 & a & b & c & d & e & f \\
		\hline
		$U/\UV$ & 0.00 & 0.50 & 1.00 & 1.50 & 2.00 & 2.50 \\
		\hline
		$I/\UA$ & 0.000 & 0.122 & 0.156 & 0.185 & 0.216 & 0.244 \\
		\hline
		\end{tabular}
		\end{center}
	\end{steps}
\end{enumerate}

%% answer: 
\jdanswer{
\begin{enumerate}
	\item
	① 1.20；② 加速度一半，0.933

	\item
	① $\times1$、8；② 如图所示；③ a、b；④ ab

\end{enumerate}
}
%% memo:
\memoanswer{
（1）① $1\Ucm + 1\Umm\times2.0$ 格 $=1.20\Ucm$；② 加速度一半，$\frac{1}{2}a=\frac{(2.8-0)\times10^{-2}}{0.06-0}\Umsq=0.467\Umsq$，所以 $a=0.933\Umsq$。

（2）① 小电珠（$2.5\UV$、$0.6\UW$）对应电阻为 $R = \frac{U^2}{P} = \frac{2.5^2}{0.6}\UO = 10.4\UO$，电阻十几欧，所以选“$\times1$”挡。② 电流表的电阻也就几欧，与电压表几千欧相比，小电珠算小电阻，所以电压表必须外接；伏安特性曲线要求电压从 0 开始逐渐变大，所以滑动变阻器必须与电压表采用分压接法。③ 从 a 端开始，电压为零，向 b 端移动，对小电珠所加电压逐渐变大。④ ab 之间电流增加了 $0.122\UA$，其它段电流增加了 $0.03\UA$ 左右，所以需要在 ab 之间将电流分为 $0.030\UA$、$0.060\UA$、$0.090\UA$，分别测出相应的电压值。

\onepicture[w=7.5cm]{figs/34_2_an.png}
}

\item
%% number: 35
%% paperName: 2011年广东理综第 35 题 0 分
%% typeId: 6
%% type: 计算题
%% chapter: 磁场
%% point: 带电粒子在复合场中的运动
%% method: 无
%% score: 0
%% degree: 350
%% duplicateId: 0
%% body:
如图（a）所示，在以 O 为圆心，内外半径分别为 $R_{1}$ 和 $R_{2}$ 的圆环区域内，存在辐射状电场和垂直纸面的匀强磁场，内外圆间的电势差 $U$ 为常量，$R_{1}=R_{0}$，$R_{2}=3R_{0}$。一电荷量为 $+q$，质量为 $m$ 的粒子从内圆上的 A 点进入，不计重力。
\begin{enumerate}
	\item
	已知粒子从外圆上以速度 $v_{1}$ 射出，求粒子在 A 点的初速度 $v_{0}$ 的大小；
	\item
	若撤去电场，如图（b）所示，已知粒子从 OA 延长线与外圆的交点 C 以速度 $v_{2}$ 射出，方向与 OA 延长线成 $45^\circ$ 角，求磁感应强度的大小及粒子在磁场中运动的时间；
	\item
	在图（b）中，若粒子从 A 点进入磁场，速度大小为 $v_{3}$，方向不确定，要使粒子一定能够从外圆射出，磁感应强度应小于多少？
\end{enumerate}

\twopicture[align=right, numA=（a）, numB=（b）, labelprefix={}, wA=4.6cm, wB=5cm]
{figs/35a.png}
{figs/35b.png}

%% answer: 
\jdanswer{
\begin{enumerate}
	\item
	$v_{0}=\sqrt{v_{1}^{2}-\frac{2qU}{m}}$

	\item
	$B=\frac{\sqrt{2}mv_{2}}{2qR_{0}}$，$t=\frac{\sqrt{2}\pi R_{0}}{2v_{2}}$

	\item
	$B'<\frac{mv_{3}}{2qR_{0}}$

\end{enumerate}
}
%% memo:
\memoanswer{
（1）根据动能定理，$qU = \frac{1}{2}mv_{1}^{2} - \frac{1}{2}mv_{0}^{2}$，所以 $v_{0} = \sqrt{v_{1}^{2} - 2qU/m}$。

（2）如图所示，设粒子在磁场中做匀速圆周运动的半径为 $R$，由几何知识可知 $R^{2} + R^{2} = (R_{2} - R_{1})^{2}$，解得 $R = \sqrt{2}R_{0}$。根据洛伦兹力公式 $qv_{2}B = m\frac{v_{2}^{2}}{R}$，解得 $B = \frac{mv_{2}}{q\sqrt{2}R_{0}} = \frac{\sqrt{2}mv_{2}}{2qR_{0}}$。

\onepicture[w=5.5cm]{figs/35_an1.png}

根据公式 $\frac{t}{T} = \frac{\theta}{2\pi}$，$2\pi R = v_{2}T$，$qv_{2}B = m\frac{v_{2}^{2}}{R}$ 解得：$t=\frac{T}{4}=\frac{2\pi m}{4Bq}=\frac{2\pi m}{4\times\frac{mv_{2}}{\sqrt{2}R_{0}}}=\frac{\sqrt{2}\pi R_{0}}{2v_{2}}$。

（3）考虑临界情况，如图所示。

\onepicture[w=5.5cm]{figs/35_an2.png}

① $qv_{3}B_{1}'=m\frac{v_{3}^{2}}{R_{0}}$，解得 $B_{1}'=\frac{mv_{3}}{qR_{0}}$；

② $qv_{3}B_{2}'=m\frac{v_{3}^{2}}{2R_{0}}$，解得 $B_{2}'=\frac{mv_{3}}{2qR_{0}}$，综合得 $B'<\frac{mv_{3}}{2qR_{0}}$。
}

\item
%% number: 36
%% paperName: 2011年广东理综第 36 题 18 分
%% typeId: 6
%% type: 计算题
%% chapter: 功和能
%% point: 动能定理
%% method: 分情况讨论
%% score: 18
%% degree: 250
%% duplicateId: 0
%% body:
如图所示，以 A、B 和 C、D 为端点的两半圆形光滑轨道固定于竖直平面内，一滑板静止在光滑的地面上，左端紧靠 B 点，上表面所在平面与两半圆分别相切于 B、C，一物块被轻放在水平匀速运动的传送带上 E 点，运动到 A 时刚好与传送带速度相同，然后经 A 沿半圆轨道滑下，再经 B 滑上滑板。滑板运动到 C 时被牢固粘连。物块可视为质点，质量为 $m$，滑板质量为 $M=2m$，两半圆半径均为 $R$，板长 $l=6.5R$，板右端到 C 的距离 $L$ 在 $R<L<5R$ 范围内取值，E 距 A 为 $S=5R$，物块与传送带、物块与滑板间的动摩擦因数均为 $\mu=0.5$，重力加速度取 $g$。
\begin{enumerate}
	\item
	求物块滑到 B 点的速度大小；
	\item
	试讨论物块从滑上滑板到离开右端的过程中，克服摩擦力做的功 $W_{f}$ 与 $L$ 的关系，并判断物块能否滑到 CD 轨道的中点。
\end{enumerate}

\onepicture[w=11cm, align=right]{figs/36.png}

%% answer: 
\jdanswer{
\begin{enumerate}
	\item
	$v_{B}=3\sqrt{gR}$

	\item
	当 $2R\leq L<5R$ 时，$W_{f}=17mgR/4$，滑块不能滑到 CD 轨道的中点；当 $R<L<2R$ 时，$W_{f}=\frac{mg(13R+2L)}{4}$。

\end{enumerate}
}
%% memo:
\memoanswer{
（1）滑块从静止开始做匀加速直线运动到 A 过程，滑动摩擦力做正功；滑块从 A 到 B，重力做正功，根据动能定理 $\mu mgS + mg\cdot2R = \frac{1}{2}mv_{B}^{2}$，解得 $v_{B} = 3\sqrt{gR}$。

（2）滑块从 B 滑上滑板后开始做匀减速运动，此时滑板开始做匀加速直线运动，当滑块与滑板达共同速度时，二者开始做匀速直线运动。设它们的共同速度为 $v$，根据动量守恒 $mv_{B}=(m+2m)v$，解得 $v=\frac{v_{B}}{3}$。

对滑块，用动能定理列方程 $-\mu mgs_{1}=\frac{1}{2}mv^{2}-\frac{1}{2}mv_{B}^{2}$，解得 $s_{1}=8R$。

对滑板，用动能定理列方程 $\mu mgs_{2}=\frac{1}{2}\times2mv^{2}-0$，解得 $s_{2}=2R$。

由此可知滑块在滑板上滑过 $s_{1} - s_{2} = 6R$ 时，小于 $6.5R$，并没有滑下去，二者就具有共同速度了。

当 $2R \leqslant L < 5R$ 时，滑块的运动是匀减速运动 $8R$，匀速运动 $L - 2R$，匀减速运动 $0.5R$，滑上 C 点，根据动能定理 $-\mu mg(8R + 0.5R) = \frac{1}{2}mv_{C}^{2} - \frac{1}{2}mv_{B}^{2}$，解得 $\frac{1}{2}mv_{C}^{2} = \frac{1}{2}mgR < mgR$，$W_{f} = \mu mg(8R + 0.5R) = 17mgR/4$，滑块不能滑到 CD 轨道的中点。

当 $R < L < 2R$ 时，滑块的运动是匀减速运动 $6.5R + L$，滑上 C 点。根据动能定理 $-\mu mg(6.5R+L)=\frac{1}{2}mv_{C}^{2}-\frac{1}{2}mv_{B}^{2}$，解得 $W_{f}=\mu mg(6.5R+L)=\frac{mg(13R+2L)}{4}$。

当 $\frac{1}{2}mv_{C}^{2}=\frac{1}{2}mg(2.5R-L)\geq mgR$ 时，可以滑到 CD 轨道的中点，此时要求 $L<0.5R$，这与题目矛盾，所以滑块不可能滑到 CD 轨道的中点。
}

\end{enumerate}

\end{document}
